\scanpage{12}
% source lines 1-4
La grande conséquence de $\uline{1^\circ}$ est la possibilité de définir $C(K)\longrightarrow G/G'$, i.e. essentiellement de savoir définir des \unsure{hom|homs}
\[
\boxed{K_{\mathfrak p}^*\longrightarrow G/G'}
\]
(notation $\left(\dfrac{x,L/K}{\mathfrak p}\right)$) de telle façon qu'on ait
\[
\boxed{\prod_{\mathfrak p}\left(\frac{x,L/K}{\mathfrak p}\right)=1\qquad\text{si }x\in K^*}
\]

% source lines 7-10
[si $\mathfrak p$ est non ramifiée, le symbole \unsure{de|d'} \deleted{Artin} \deleted{est} \unsure{canonique} étant défini par l'hom de \unsure{Frobénius|Frobènius|Frobênius}]. Le lien avec la \uline{théorie des corps de classes \unsure{local|locale}} est le suivant :

% source lines 11-15
On a pour tout $\mathfrak p$, \unsure{des hom naturels|des homs naturels|un hom naturel}
\diagram{\begin{tikzcd}[ampersand replacement=\&]
A(K_{\mathfrak p})\arrow[r,"\text{surj.}"]\&
\begin{gathered}\text{groupe de décomposition }G_{\mathfrak p}^{d}\\\text{de }G\text{ relatif à }\mathfrak p\end{gathered}
\end{tikzcd}}
(on suppose $L/K$ abélien pour simplifier) où $A(K_{\mathfrak p})$ désigne le groupe de \unsure{galois|Galois} de l'extension abélienne maximale de $K_{\mathfrak p}$. D'autre part d'après la théorie du corps de classes \unsure{local|locale|locales} on a un isomorphisme canonique
% archive sheet number: 11
